\subsection{Проверка знаний}
Перед выполнением практических заданий ответьте на вопросы. Вставьте пропущенную строчку кода или слово. Выберите один верный вариант и сверьте с пояснением.
\begin{enumerate}
\item Какое значение размера буфера должно быть передано в инициализацию? \verb|leds = Ledbar.new(___)|
  \begin{enumerate}[label=(\alph*)]
    \item 29 \item 8 \item 4 \item 1
  \end{enumerate}
  Ответ: (a). Пояснение: драйвер считывает буфер из 29 элементов.

  \item Чем должно быть начальное значение счётчика в цикле? \verb|for i = ____, ledNumber-1 do ... end|
  \begin{enumerate}[label=(\alph*)]
    \item 0 \item 1 \item -1 \item ledNumber
  \end{enumerate}
  Ответ: (a). Пояснение: индексы линейки начинаются с 0 и заканчиваются 28.

  \item Каким вызовом передать RGB-таблицу в \verb|leds:set| при объявленном \verb|local unpack = table.unpack|?
  \begin{enumerate}[label=(\alph*)]
    \item \verb|leds:set(i, col)| \item \verb|leds:set(i, unpack(col))| \item \verb|leds:set(i, toRGB(col))| \item \verb|leds:set(i, rgb(col))|
  \end{enumerate}
  Ответ: (b). Пояснение: \verb|unpack| распаковывает таблицу \verb|{r,g,b}| в три аргумента.

  \item Как корректно единожды включить зелёный цвет на всей линейке?
  \begin{enumerate}[label=(\alph*)]
    \item \verb|changeColor({0,1,0})| \item \verb|changeColor({1,0,0})| \item \verb|changeColor({0,0,1})| \item \verb|changeColor({1,1,1})|
  \end{enumerate}
  Ответ: (a). Пояснение: компоненты RGB для зелёного — \verb|{0,1,0}|.

  \item Как создать отложенный запуск синего через 3 секунды?
  \begin{enumerate}[label=(\alph*)]
    \item \verb|sleep(3); changeColor({0,0,1})| \item \verb|Timer.new(3, changeColor)| \item \verb|Timer.callLater(3, function() changeColor({0,0,1}) end)| \item \verb|ap.push(Ev.COPTER_LANDED)|
  \end{enumerate}
  Ответ: (c). Пояснение: \verb|Timer.callLater| планирует единичный вызов без блокировки.

  \item Какой период задать для мигания раз в полсекунды? \verb|Timer.new(___, fn)|
  \begin{enumerate}[label=(\alph*)]
    \item 0.5 \item 500 \item 50 \item 1.5
  \end{enumerate}
  Ответ: (a). Пояснение: период указывается в секундах.

  \item Как инициализировать булевое состояние для мигания?
  \begin{enumerate}[label=(\alph*)]
    \item \verb|on = false| \item \verb|on = 0| \item \verb|on = "false"| \item \verb|on = nil|
  \end{enumerate}
  Ответ: (a). Пояснение: булев тип Lua использует \verb|true/false|.

  \item Как проверить чётность индекса \verb|i|?
  \begin{enumerate}[label=(\alph*)]
    \item \verb|i / 2 == 0| \item \verb|i % 2 == 0| \item \verb|i == even| \item \verb|isEven(i)|
  \end{enumerate}
  Ответ: (b). Пояснение: оператор остатка \verb|%| определяет чётность.

  \item Как остановить запущенный таймер \verb|t|?
  \begin{enumerate}[label=(\alph*)]
    \item \verb|t:stop()| \item \verb|t:pause()| \item \verb|stop(t)| \item \verb|t.cancel()|
  \end{enumerate}
  Ответ: (a). Пояснение: API таймера предоставляет методы \verb|start/stop|.

  \item Как запустить таймер \verb|t| после создания?
  \begin{enumerate}[label=(\alph*)]
    \item \verb|t:run()| \item \verb|t:start()| \item \verb|start(t)| \item \verb|t:on()|
  \end{enumerate}
  Ответ: (b). Пояснение: корректный метод — \verb|start|.

  \item Как называется функция-обработчик событий автопилота?
  \begin{enumerate}[label=(\alph*)]
    \item \verb|function events(e)| \item \verb|function callback(event)| \item \verb|function onEvent(e)| \item \verb|function handler(ev)|
  \end{enumerate}
  Ответ: (b). Пояснение: платформа вызывает \verb|callback| и передаёт \verb|event|.

  \item Какое событие соответствует посадке?
  \begin{enumerate}[label=(\alph*)]
    \item \verb|Ev.POINT_REACHED| \item \verb|Ev.ALTITUDE_REACHED| \item \verb|Ev.COPTER_LANDED| \item \verb|Ev.LOW_VOLTAGE2|
  \end{enumerate}
  Ответ: (c). Пояснение: \verb|COPTER_LANDED| — системное событие посадки.

  \item Какое событие соответствует низкому напряжению?
  \begin{enumerate}[label=(\alph*)]
    \item \verb|Ev.LOW_VOLTAGE2| \item \verb|Ev.BATTERY_OK| \item \verb|Ev.POWER_ON| \item \verb|Ev.MOTOR_ARMED|
  \end{enumerate}
  Ответ: (a). Пояснение: используйте \verb|LOW_VOLTAGE2| для аварийной индикации.

  \item Как выглядит условие для «полосы загрузки»: включить все индексы \verb|<= k|?
  \begin{enumerate}[label=(\alph*)]
    \item \verb|if i < k then ...| \item \verb|if i <= k then ...| \item \verb|if k <= i then ...| \item \verb|if i == k then ...|
  \end{enumerate}
  Ответ: (b). Пояснение: включёнными остаются все индексы не выше текущего \verb|k|.

  \item Какой допустимый диапазон значений \verb|r,g,b| в модели RGB?
  \begin{enumerate}[label=(\alph*)]
    \item \verb|0..255| \item \verb|0..100| \item \verb|0.0..1.0| \item \verb|-1..1|
  \end{enumerate}
  Ответ: (c). Пояснение: платформа принимает вещественные значения от \verb|0.0| до \verb|1.0|.
\end{enumerate}
